\chapter*{Introduction}
    \section*{About This Document}
    {\color{red}TODO}    

    \section*{A Fairytale}
    At many places in homotopy theory interesting maps are uniquely determined up to homotopy but not on the nose. Let me give 4 examples:

    \begin{enumerate}[label=\arabic*)]
        \item 
            The inverse to a homotopy equivalence \label{list01-01}
        \item 
            The classifying map $B \to BO(d)$ of a $ d$-dimensional vector bundle \label{list01-02}
        \item 
            The lift in a HELP-diagram with one of the vertical maps a homotopy equivalence \label{list01-03}
            \begin{equation*}
                \begin{tikzcd}[row sep=scriptsize, column sep=scriptsize]
                    A && E \\
                    \\
                    X && B
                    \arrow[from=1-1, to=1-3]
                    \arrow["{\color{magenta}{\text{cofibration} \longrightarrow}}"', hook, from=1-1, to=3-1]
                    \arrow["{\color{magenta}{\longleftarrow \text{fibration}}}", from=1-3, to=3-3]
                    \arrow[dashed, from=3-1, to=1-3]
                    \arrow[from=3-1, to=3-3]
                \end{tikzcd}
            \end{equation*} 
        \item 
            The $n$-fold multiplication
            \begin{equation*}
                \left( \Omega X \right)^n \to \Omega X
            \end{equation*}\label{list01-04}
        Similar phenomena occur in algebra:
        \item 
            The map $F \to F'$ between free resolutions of $M \to N$. \label{list01-05}
    \end{enumerate}

    Now, faced with such a situation it is often difficult to perform categorically flavoured operations in a meaningful way:

    A simple example is the formation of fibres: $f^{-1}$ is the pullback
    \begin{equation*}
        \begin{tikzcd}
            f^{-1}(y) \arrow[r] \arrow[d]
            \arrow[dr, phantom, "\lrcorner", very near end]
            & X \arrow["f", d] \\
            y \arrow["i"', r]
            & Y 
        \end{tikzcd}        
    \end{equation*}
    but homotopic $f$ and $f'$ (or homotopic $i$, $i'$) need not give homotopy equivalent pullbacks.
    To assign invariant meaning to them homotopical versions of such constructions were soon invented, in the case at hand, the \emph{homotopy fibre}.

    But this still not quite good enough:
    \begin{enumerate}[label=\arabic*)]
        \item While homotopy fibres of homotopic maps $f$ and $f'$ are homotopy equivalent, the (homotopy class of) homotopy equivalence depends on a choice of (homotopy class of) homotopy between $f$ and $f'$, so the homotopy fibre is not functorial on the arrow category of the homotopy category.
        \item The homotopy fibre is not a direct categorical construction anymore neither in the original category nor the homotopy category; it represents neither the functor of commuting diagrams nor of homotopy commuting diagrams. Rather a map 
        \begin{equation*}
            Z \to \hofib_x (f)
        \end{equation*}
        represents a \emph{homotopy coherent} diagram
        \begin{equation*}
            \begin{tikzcd}[column sep=scriptsize, row sep=scriptsize]
                Z && X \\
                & {} \\
                * && Y
                \arrow[color={rgb,255:red,255;green,0;blue,0}, from=1-1, to=1-3]
                \arrow[color={rgb,255:red,255;green,0;blue,0}, from=1-1, to=3-1]
                \arrow[color={rgb,255:red,255;green,0;blue,0}, from=1-1, to=3-3]
                \arrow[color={rgb,255:red,255;green,0;blue,0}, Rightarrow, from=1-3, to=2-2]
                \arrow["f", from=1-3, to=3-3]
                \arrow[color={rgb,255:red,255;green,0;blue,0}, Rightarrow, from=3-1, to=2-2]
                \arrow["x"', from=3-1, to=3-3]
            \end{tikzcd}
        \end{equation*}
        (double arrows indicating homotopies)
    \end{enumerate}

    The first problem is addressed by the observation that the maps in question are all determined up to contractible choice not just up to homotopy! Let me warn you that it is not true in examples \ref{list01-01}, \ref{list01-02}, \ref{list01-04} and \ref{list01-05} that the component of a given map in the mapping space is contractible, just that there is a contractible "space", in which the maps in question naturally live: for example in \ref{list01-02} the space of bundle maps
    \begin{equation*}
        V \to EO(d) \times_{O(d)} \mathbb{R}^d
    \end{equation*}
    is contractible, in \ref{list01-04} the multiplications are given partitions of the interval into 4 parts i.e. the space
    \begin{equation*}
        \left \{ (t_1, t_2, t_3) \in \II^3 \mid 0 < t_1 < t_2 < t_3 < 1 \right \}
    \end{equation*}
    parametrizes the multiplications and is contractible. For a homotopy equivalence $g \, \colon \, Y \to X$ it is a space of pairs
    \begin{equation*}
        \left \{ (f, H) \in C(X, Y) \times C(X \times \II, X) \mid H_0 = \id_X, H_1 = g \circ f \right \} \,.
    \end{equation*}

    In particular, for any two choices of map as in \ref{list01-01}-\ref{list01-05} not only does there exist a homotopy between them, but that homotopy is itself unique up to homotopy (so for example for a vector bundle $V \to B$ the homotopy fibre
    \begin{equation*}
        \hofib_{\delta}(B \to BSO(d))
    \end{equation*}
    is well-defined up to unique isomorphism in the homotopy category (once a point $\delta \in BSO(d)$ is fixed once and for all). But moreover any two homotopies witnessing this uniqueness are again unique up to homotopy, so further constructions of similar flavour are again possible without destroying functoriality.
    To exploit these observations a systematic category theory allowing for homotopy and contractibility instead of equality and uniqueness is required. The easiest way to achieve this is by considering topologically enriched categories. To accommodate non-topological examples it is in fact better to work combinatorially and consider Kan-enriched categories. While the entirety of higher category theory can be developed for these objects, they carry a major defect: composition of morphisms is still strictly associative. While this is the case for most categories that serve as starting points ($\Top$, $\sSet$, $\Ch_R$, ...), it is not preserved by many categorical constructions, since (by design) these are only determined up to contractible choice.
    A related problem is that a functor $F \colon \Cc \to \Dd$ of simplicial categories ought to be considered invertible if
    \begin{enumerate}[label=\arabic*)]
        \item it is essentially homotopically surjective, meaning for all $D \in \Dd$ there is a $C \in \Cc$ and a homotopy equivalence $F(C) \simeq D$ and
        \item $F$ induces a homotopy equivalence 
        \begin{equation*}
            \Hom(A, B) \to \Hom(FA, FB)
        \end{equation*}
        for all objects $A, B \in \Cc$.
    \end{enumerate}
    Such functors are called \emph{Dwyer-Kan equivalences}. 

    Although this notion restricts to usual equivalences of categories for discretely enriched $\Cc$ and $\Dd$ it is quite different from that of an enriched equivalence, which would require isomorphisms 
    \begin{equation*}
        \Hom(A, B) \to \Hom(FA, FB) \,.
    \end{equation*}
    Here are two examples. Consider the full subcategory $\Cc$ of the simplicially enriched category $\Top$ spanned by contractible spaces. Then the functor $p \colon \Cc \to *$ to the one-point category (with trivial enrichment) is a Dwyer-Kan equivalence, since the space of continuous functions between contractible spaces is itself contractible. However, for any functor $s \colon * \to \Cc$ a natural transformation $\tau \colon s \circ p \Rightarrow \id_{\Cc}$ would require all diagrams of the form
    \begin{equation*}
        \begin{tikzcd}
            s(*) \arrow["\tau_X"', r] \arrow["(s \circ p)(f)"', d]
            & X \arrow["f", d] \\
            s(*) \arrow["\tau_Y"', r]
            & Y 
        \end{tikzcd}        
    \end{equation*}
    to commute, which is clearly absurd since $s(*) \neq \varnothing$. Note, however, that a (unique) transformation $\id_{\Cc} \Rightarrow s \circ p$ exists just fine, showing that natural transformations, which are pointwise equivalences need not admit inverses, another desirable criterion that Kan-enriched categories fit. As a second example consider the category $\Nn[2]$ with three objects and six morphisms
    \begin{equation*}
        \begin{tikzcd}[column sep=small]
                                                                                       & 1 \arrow[rd, "g"] \arrow["\id"', loop, distance=2em, in=125, out=55] &                                                      \\
0 \arrow[ru, "f"] \arrow[rr, "gf"'] \arrow["\id"', loop, distance=2em, in=215, out=145] &                                                                      & 2 \arrow["\id"', loop, distance=2em, in=35, out=325]
\end{tikzcd}
    \end{equation*}
    We can give it the trivial enrichment.
    Consider also the simplicially enriched category $\Cc$, whose only difference is that we replace 
    \begin{equation*}
        \Hom_{\Nn[2]}(0, 2) = \const \{f, g\} \in \sSet
    \end{equation*}
    by 
    \begin{equation*}
        \Hom_{\Cc}(0, 2) = \Sing \II
    \end{equation*}
    with $g \circ f = (\Delta^0 \overset{0}{\longrightarrow} \II) \in \Sing_0 \II$.
    Evidently, there are unique enriched functors
    \begin{equation*}
        \Cc \to \Nn[2] \quad \text{and} \quad \Nn[2] \to \Cc
    \end{equation*}
    which are identity on objects and both are Dwyer-Kan equivalences. However, only one make the diagram
    \begin{equation*}
        \begin{tikzcd}
                                                                      &                                                            &                                   &                                                                                                                                                                                                                                                   & {\mathcal{N}[2]} \arrow[dd, "\sim"', shift right] \\
            \overset{0}{\bullet} \arrow["\id"', loop, distance=2em, in=215, out=145] \arrow[r] & \overset{1}{\bullet} \arrow["\id"', loop, distance=2em, in=35, out=325] & {} \arrow[r, equal] & {\mathcal{N}[1]} \arrow[rd, red, maps to, shift right=15, swap, , "(0 \to 1)" at start, "\Delta^0 \overset{1}{\to} \mathrm{I}" at end] \arrow[ru, red, maps to, shift left=9, "0" at start, "0" at end] \arrow[ru, red, maps to, shift left=4, "1" at start, "2" at end] \arrow[ru, hook] \arrow[rd, red, maps to, shift right=11, , "1" at start, "2" at end] \arrow[rd, red, maps to, shift right=6, "0" at start, "0" at end] \arrow[rd, hook, shift right] &                                                   \\
                                                                      &                                                            &                                   &                                                                                                                                                                                                                                                   & \Cc \arrow[uu, "\backsim"', shift right=2]       
        \end{tikzcd}
    \end{equation*}
    commute. The worst offender, however, are Dwyer-Kan equivalences which do not admit inverses at all: for example, there is no reason to expect that the Dwyer-Kan equivalence
    \begin{equation*}
        \left\{   
            \begin{array}{c} 
                \text{finite cell complexes,} \\
                \text{cellular\footnotemark maps}
            \end{array}
        \right\} 
        \subseteq 
        \left\{
            \begin{array}{c}
                \text{spaces homotopy equivalent to a finite cell complex,} \\
                \text{continuous maps}
            \end{array}
        \right\}
    \end{equation*}
    \footnotetext{Suitably interpreted as a simplicial set the statement fails for the naive subspace of cellular maps inside all maps!}
    admits any functor the other way that induces the inverse on homotopy categories.
    To remedy these deficits one considers a notion of higher category theory that is not equipped with a strict composition law. There are at least two good candidates: Joyal's \emph{quasi-categories} and Rezk's \emph{complete Segal spaces}. Given the recent development of the former at the hands of Lurie we will give preference to it and say little about the later notion.
    Quasi-categories were originally defined by Boardman and Vogt (who taught me topology in Osnabrück) during the development of $E_{\infty}$-spaces (the natural notion of commutative monoid that arises in higher categories): they attempted to define the morphisms between such objects as maps of the underlying spaces together with an $E_{\infty}$-structure on its cofibre extending the given ones.
    As cofibres do not commute with composition they did not quite obtain a category and instead formalized the resulting object (under the name "weak Kan complex").
    This apparent defect was soon remedied by May's reformulation of their theory via his introduction of operads and thus stayed undeveloped for many years until Joyal demonstrated that much of ordinary category theory generalises to such weak Kan complexes (whence he renamed them). A quasi-category is (roughly speaking) a simplicial set $\Cc$, which admits compositions up to contractible choice; here one thinks of $\Cc_0$ as the objects and $\Cc_1$ as the morphisms of our higher category. The space of composites of two 1-simplices $f, g \in \Cc_1$ is then the pullback
    \begin{equation*}
        \begin{tikzcd}
            \operatorname{Comp}(f,g) \arrow[r] \arrow[d]
            \arrow[dr, phantom, "\lrcorner", very near end]
            & \Hom(\Delta^2, \Cc) \arrow[d] \\
            * \arrow[r, "{(f,\; g)}"]
            & \Hom(\Lambda^2_1, \Cc) 
        \end{tikzcd}        
    \end{equation*}    
    \emph{Warning}: the actual definition requires the right hand vertical map to also be a Kan fibration!
    
    We will develop the very basics of this notion. The main results we will hopefully show, say precisely that quasi-categories do not share the defects of Kan-enriched categories, i.e. that
    \begin{enumerate}[label=\arabic*)]
        \item functors are equivalences if they are fully faithful and essentially surjective (both in the homotopical sense) and
        \item natural transformations are invertible if they consist of equivalences.
    \end{enumerate}
    Furthermore, we will show that the classical homotopy theory of cell complexes is precisely the category theory of higher groupoids, a fact known as Grothendieck's \emph{homotopy hypothesis}.
    
    By way of compensation it is, however, quite a bit more difficult to produce interesting quasi-categories directly. We will therefore develop the \emph{homotopy coherent nerve}, which produces a quasi-category from a Kan-enriched ordinary category. Both these last points require a bit a classical simplicial homotopy theory (which is the theory of Kan complexes) so we will treat that first.
    
    Finally, let me just mention that the theories just discussed are only the start: they are all on models for $(\infty, 1)$-category theory, meaning that all morphisms of degree $\geq 2$ (starting with homotopies) are automatically invertible. For example, the (large) $(\infty, 1)$-category $\operatorname{Cat}_{\infty}$ of (small) $(\infty, 1)$-categories really wants to be an $(\infty, 2)$-category (since not all natural transformations are equivalences).
    
    We will have little to say on this. Other topics that will not make an appearance here:
    \begin{itemize}[label=$-$]
        \item straightening / unstraightening
        \item Yoneda's Lemma
        \item operads / symmetric monoidal categories
        \item sheaves / hypersheaves
        \item spectra and stabilization
        \item systematic study of limits and colimits
    \end{itemize}

    Maybe if there is a part 2...